% Part: normal-modal-logic
% Chapter: syntax-and-semantics

\documentclass[../../../include/open-logic-chapter]{subfiles}

\begin{document}

\olchapter{nml}{syn}{语法与语义}

模态逻辑通过添加额外的算子来扩展经典逻辑，旨在形式化“必然性”与“可能性”这类概念。例如，“必然为真”和“可能为真”就是这样的概念。虽然经典逻辑的联结词可以组合起来表达复杂的命题，但它们无法表达关于命题以何种方式为真的命题。经典逻辑没有提供任何表达式，使我们能够表达一个真命题是必然的还是偶然的，而不仅仅断言它是真的或假的。因此，为了表达这些想法，我们通过引入新的算子来扩展其语言。

\olimport{introduction}
\olimport{language-modal-logic}
\olimport{substitution}
\olimport{relational-models}
\olimport{truth-at-w}
\olimport{truth-in-model}
\olimport{modal-validity}
\olimport{tautological-instances}
\olimport{schemas}
\olimport{entailment}

\OLEndChapterHook

\end{document}